%==============================================================================
% CHAPTER 11: The Mathematics of Cyber Defense
%==============================================================================
\chapter{The Mathematics of Cyber Defense}
\label{ch:cyber} \cite{emanuello2022cyber}

\begin{flushright}
\textit{John A. Emanuello, Ahmad Ridley}\\
Communicated by \textit{Notices} Associate Editor Emilie Purvine
\end{flushright}

\epigraph{``The speed, complexity, and ubiquity of cyber-attacks has never been more apparent... Current cyber-defense capabilities are static and rules-based... This approach is unsustainable.''}{--- Emanuello \& Ridley} \cite{emanuello2022cyber}

\begin{abstract}
Cyber defense requires detecting anomalies in massive, heterogeneous data streams (logs, netflow, alerts) and responding autonomously. This chapter surveys the mathematical foundations: autoencoder\index{autoencoders}s for unsupervised anomaly detection, \index{anomaly detection}Word2Vec-style embeddings for categorical log data, and reinforcement learning \index{reinforcement learning, cybersecurity}for autonomous response. The unique challenges: adversarial dynamics, non-stationary distributions, and the need for human-machine teaming, \cite{emanuello2022cyber}.
\end{abstract}

%==============================================================================
\section{Introduction: Why Mathematics for Cyber?}

\noindent\textit{Mathematical notation follows standard conventions. Key terms are defined in the glossary.}\n
\label{sec:cyber:intro} \cite{emanuello2022cyber}

Cyber attacks are not single events but \emph{campaigns}: multi-phase, adaptive, stealthy. Defenders face:
\begin{itemize}
    \item \textbf{Volume}: Terabytes/day of logs, netflow, packet captures
    \item \textbf{Heterogeneity}: Categorical (IPs, usernames), temporal, graph-structured
    \item \textbf{Adversariality}: Attackers adapt to defenses
    \item \textbf{Label scarcity}: Almost no labeled attacks; must detect \emph{anomalies}
\end{itemize}

Mathematical approach: \textbf{Build models of normal behavior; flag deviations.}

%==============================================================================
\section{Data Representation: Embedding Categorical Features}
\label{sec:cyber:embedding} \cite{emanuello2022cyber}

Log data is nominal: IP addresses, usernames, process names, file paths. Standard ML needs numeric vectors.

\textbf{Word2Vec for logs (Mikolov et al., 2013)}: Treat a log line as a ``sentence,'' tokens as ``words.'' Train skip-gram model:
\[
\max_{\theta} \sum_{(w,c) \in D} \log P(c|w;\theta), \quad P(c|w) = \frac{\exp(\bm{v}_c^\top \bm{v}_w)}{\sum_{c'} \exp(\bm{v}_{c'}^\top \bm{v}_w)}.
\]

Result: Each token (IP, process, etc.) gets a dense vector $\bm{v} \in \mathbb{R}^d$. Similar behavior $\to$ nearby vectors.

\begin{lstlisting}[style=python,caption=Log embedding with Word2Vec]
from gensim.models import Word2Vec
import numpy as np

# Parse logs into token sequences
log_sequences = [
    ["192.168.1.1", "ssh", "accept", "user1"],
    ["192.168.1.2", "ssh", "failed", "user2"],
    # ...
]

# Train Word2Vec
model = Word2Vec(log_sequences, vector_size=100, window=5, min_count=1, workers=4)

# Embed a log line
def embed_log(tokens, model):
    vecs = [model.wv[t] for t in tokens if t in model.wv]
    return np.mean(vecs, axis=0) if vecs else np.zeros(model.vector_size)

# Now each log is a 100-d vector for downstream ML
\end{lstlisting}

%==============================================================================
\section{Anomaly Detection: Deep Autoencoders}
\label{sec:cyber:ae} \cite{emanuello2022cyber}

\textbf{Deep Autoencoder (AE)}: Multi-layer perceptron trained to reconstruct input:
\[
\bm{z} = \sigma(\bm{W}_e \bm{x} + \bm{b}_e), \quad \hat{\bm{x}} = \sigma(\bm{W}_d \bm{z} + \bm{b}_d),
\]
loss $= \|\bm{x} - \hat{\bm{x}}\|^2$. The bottleneck $\bm{z}$ learns a compressed representation.

\textbf{Unsupervised anomaly detection}: Train on \emph{normal} traffic only. At test time, high reconstruction error $\to$ anomaly, \cite{emanuello2022cyber}.

\begin{theorem}[AE as Nonlinear PCA]
A deep AE with linear activations learns the PCA subspace. With nonlinearities, it learns a \emph{nonlinear manifold} of normal data.
\end{theorem}

\begin{lstlisting}[style=python,caption=Autoencoder for log anomaly detection] \cite{emanuello2022cyber}
import torch
import torch.nn as nn

class LogAutoencoder(nn.Module):
    def __init__(self, input_dim, latent_dim=32):
        super().__init__()
        self.encoder = nn.Sequential(
            nn.Linear(input_dim, 128), nn.ReLU(),
            nn.Linear(128, 64), nn.ReLU(),
            nn.Linear(64, latent_dim)
        )
        self.decoder = nn.Sequential(
            nn.Linear(latent_dim, 64), nn.ReLU(),
            nn.Linear(64, 128), nn.ReLU(),
            nn.Linear(128, input_dim)
        )
    
    def forward(self, x):
        z = self.encoder(x)
        return self.decoder(z), z

# Training on normal data only
model = LogAutoencoder(input_dim=100)
optimizer = torch.optim.Adam(model.parameters(), lr=1e-3)

for epoch in range(50):
    for batch in normal_dataloader:
        recon, _ = model(batch)
        loss = nn.MSELoss()(recon, batch)
        loss.backward()
        optimizer.step()
        optimizer.zero_grad()

# Anomaly score = reconstruction error
def anomaly_score(model, x):
    with torch.no_grad():
        recon, _ = model(x)
        return torch.mean((x - recon)**2, dim=1)
\end{lstlisting}

%==============================================================================
\section{Chaining Anomalies Across Modalities}
\label{sec:cyber:chaining} \cite{emanuello2022cyber}

A single anomaly (e.g., unusual process) may be benign. A \emph{chain} across modalities (unusual process + unusual netflow + unusual file access) is likely an attack.

\textbf{Challenge}: Fuse AE scores from log embeddings, netflow embeddings, alert embeddings into a unified timeline.

Approach: \textbf{Temporal point processes} or \textbf{graph neural networks} on the alert correlation graph.

%==============================================================================
\section{Autonomous Response: Reinforcement Learning}
\label{sec:cyber:rl} \cite{emanuello2022cyber}

Detection is not enough; must \emph{respond}. Formulate as \textbf{Markov Decision Process}:
\begin{itemize}
    \item State $s_t$: Network state, active alerts, system health
    \item Action $a_t$: Block IP, isolate host, rate-limit, collect forensics
    \item Reward $r_t$: Negative for damage, positive for containment, cost for disruption
\end{itemize}

Learn policy $\pi(a|s)$ via \textbf{Deep Q-Network (DQN)} or \textbf{Proximal Policy Optimization (PPO)}.

\begin{lstlisting}[style=python,caption=RL environment for cyber response] \cite{emanuello2022cyber}
import gymnasium as gym
import numpy as np

class CyberEnv(gym.Env):
    def __init__(self):
        # State: [num_alerts, severity_scores, network_health, ...]
        self.observation_space = gym.spaces.Box(-np.inf, np.inf, (20,))
        # Actions: 0=monitor, 1=block_ip, 2=isolate_host, 3=rate_limit, 4=forensics
        self.action_space = gym.spaces.Discrete(5)
        self.state = None
        self.attack_ongoing = False
    
    def step(self, action):
        # Simulate attack progression and response effect
        reward = self._compute_reward(action)
        terminated = self.attack_neutralized()
        truncated = self.steps >= 100
        return self.state, reward, terminated, truncated, {}
    
    def reset(self, seed=None):
        self.state = np.zeros(20)
        self.attack_ongoing = np.random.rand() < 0.3
        return self.state, {}
\end{lstlisting}

\textbf{Human-in-the-loop}: RL agent proposes actions; human analyst approves/modifies. Over time, trust increases $\to$ more autonomy.

%==============================================================================
\section{Bridge to Chapter 12}
\label{sec:cyber:bridge} \cite{emanuello2022cyber}

Cyber defense uses deep learning for \emph{pattern recognition} in dynamic systems. Chapter 12 uses \textbf{operads} (categorical algebra) for \emph{compositional design} of complex systems. Both address the challenge of \emph{structure in complexity}: one learns it, the other specifies it.

%==============================================================================
\section{Further Reading}
\label{sec:cyber:refs} \cite{emanuello2022cyber}

\begin{itemize}
    \item Emanuello, Ridley (2022). \textit{The Mathematics of Cyber Defense}. Notices, 69(6).
    \item Golczynski, Emanuello (2021). \textit{End-to-End Anomaly Detection...}. arXiv:2108.12276.
    \item Sutton, Barto (2018). \textit{Reinforcement Learning}. MIT Press.
    \item Goodfellow et al. (2016). \textit{Deep Learning}. MIT Press (Ch. 14 on AEs).
    \item MITRE ATT\&CK framework: \url{https://attack.mitre.org}
\end{itemize}

%==============================================================================
\section{Exercises}
%==============================================================================
% Solutions
%==============================================================================
\section{Solutions}
\label{sec:cyber:solutions}

\begin{solution}
\textbf{Exercise 1 (Autoencoder):} An autoencoder with hidden dimension $k < n$ learns to compress log data. The reconstruction error $\|x - \hat{x}\|^2$ identifies anomalies: large errors indicate unusual patterns.
\end{solution}

\begin{solution}
\textbf{Exercise 2 (Word2Vec):} Word2Vec learns embeddings for categorical data (IP addresses, usernames). The cosine similarity between embeddings captures semantic relationships: similar IPs have similar usage patterns.
\end{solution}

\begin{solution}
\textbf{Exercise 3 (RL for defense):} A reinforcement learning agent learns to respond to attacks by maximizing a reward function. The state space includes network metrics; the action space includes firewall rules and alerts.
\end{solution}

\begin{solution}
\textbf{Exercise 4 (Research):} Adversarial attacks on anomaly detection systems can fool autoencoders by adding small perturbations. Defense mechanisms include adversarial training and ensemble methods.
\end{solution}

\label{sec:cyber:exercises} \cite{emanuello2022cyber}

\begin{exercise}
Train a Word2Vec model on the HDFS log dataset. Visualize the embedding of IP addresses---do malicious IPs cluster?
\end{exercise}

\begin{exercise}
Implement a variational autoencoder (VAE) for log anomalies. Compare reconstruction error vs. ELBO as anomaly scores.
\end{exercise}

\begin{exercise}
Design a reward function for a cyber RL agent that balances: (1) minimizing data exfiltration, (2) minimizing false positive isolations, (3) minimizing analyst workload, \cite{emanuello2022cyber}.
\end{exercise}

\begin{exercise}[Research]
Investigate \emph{adversarial attacks on the anomaly detector itself}: can an attacker craft traffic that looks normal to the AE but is malicious? \cite{emanuello2022cyber}
\end{exercise}